\documentclass[10pt,a4paper]{amsart}
\usepackage[margin=19mm]{geometry}
\usepackage{amsmath,amssymb,mathtools}
\usepackage[scr=rsfs,cal=euler]{mathalfa}
\usepackage{xfrac}
\usepackage[unicode,hidelinks,bookmarks=false]{hyperref}
\newcommand{\ie}{i.e.\ }
\newcommand{\cf}{cf.\ }
\newcommand{\resp}{respectively\ }
\newcommand{\Subsection}{\subsection}
\providecommand{\nobreakdash}{\nobreak\hbox{-}}
\setlength{\emergencystretch}{2em}
\multlinegap0pt
\usepackage{mathtools}
\usepackage{amssymb}
\usepackage[shortlabels]{enumitem}
\usepackage{tikz}
\usepackage{xcolor}
\newtheorem{lemma}{Lemma}
\newcommand{\dd}{\mathrm{d}}
\title{On the spectral radius of the ratio of Girko matrices}
\author{Djalil Chafaï, David García-Zelada, Yuan Yuan Xu}
\date{Source: arXiv:2510.18669v3 (CC BY 4.0)}
\begin{document}
\maketitle

\noindent\emph{Source and licence.} On the spectral radius of the ratio of Girko matrices. Version arXiv:2510.18669v3 (CC BY 4.0), distributed by arXiv under the Creative Commons Attribution 4.0 International licence, http://creativecommons.org/licenses/by/4.0/, as recorded in the arXiv licence metadata for this exact version and shown on the abstract page https://arxiv.org/abs/2510.18669v3. Full licence notice: \texttt{licenses/arxiv-2510.18669-CC-BY-4.0.txt}.\par\smallskip

\noindent\textit{Editorial scope.} Counted unit: the article's lemma lem:ChangeOfCoordinates (source0.tex lines 1539--1558). The article's complete original proof of it is delivered with it (source0.tex lines 1560--1631, one \texttt{proof} environment of 72 lines). No mathematics is written, repaired or re-derived: the statement and the proof are the article's own lines, in the article's own notation, and no retained line is substituted or re-flowed.  No further source-local context is needed: every cross-reference of every retained line resolves inside the two delivered spans.  Nothing was omitted from any retained span, so the package carries no omission record.  No retained line contains a citation, so the wrapper carries no bibliography and no bibliography entry.  No retained line contains a figure environment, an image-inclusion command or drawing code, so no image is delivered and the package declares no asset dependency.  Editorial formatting only, in addition: an AMS \texttt{amsart} preamble that restates the article's own declarations and macros used by the retained lines, namely the environments \texttt{lemma} on separate counters, together with the standard packages those lines need.  The retained environments print in this document as Lemma 1 in order of appearance, whereas the article prints the same items with the numbering of its own full document.  The complete native source of the article as distributed in the arXiv e-print for this exact version is bundled unchanged as \texttt{sources/187517/source0.tex}.
\begin{lemma}[Delta-method for point processes] \label{lem:ChangeOfCoordinates}
Let $(\mathcal X_n)_{n \geq 1}$ be a sequence
of point processes on an open set 
$U \subset \mathbb C$ and fix $x_0 \in U$. Consider two open sets $V_1,V_2 \subset \mathbb C$
containing the origin 
together with diffeomorphisms
$\varphi_i: U \to V_i$.
 Then
$\sqrt n (\varphi_1(\mathcal X_n) - 
\varphi_1(x_0))$
converges in law if and only if $\sqrt n (\varphi_2(\mathcal X_n) - \varphi_2(x_0))$
also does and, in that case,
\[\lim_{n \to \infty}
\sqrt n \big(\varphi_1(\mathcal X_n) - \varphi_1(x_0) \big)
=
(\mathrm d \varphi_1)_{x_0} \circ
(\mathrm d \varphi_2)_{x_0}^{-1}
\Big(\lim_{n \to \infty}
\sqrt n \big(\varphi_2(\mathcal X_n) - \varphi_2(x_0)\big)\Big).\]
\end{lemma}

\begin{proof}
First, by changing $\varphi_i$ by 
$\varphi_i - \varphi_i(x_0)$, we may consider $\varphi_i(x_0)=0$.
Then, under this assumption, by considering the point process 
$\mathcal Y_n = \varphi_1(\mathcal X_n)$ 
and the maps
$\psi_i = \varphi_i \circ \varphi_1^{-1}$
the statement reduces to the following statement.
Let $\mathcal Y_n$ be a sequence of point
processes on $V_1$ that contains the origin 
such that 
$\sqrt n \mathcal Y_n$ converges. If
$\psi: V_1 \to V_2$ is a diffeomorphism
that satisfies $\psi(0) = 0$ then
\[\lim_{n \to \infty} 
\sqrt n\psi(\mathcal Y_n)
=(\mathrm d \psi)_0\big(\lim_{n \to \infty} \sqrt n \mathcal Y_n \big) .\]
To show this, let us consider 
a compactly supported smooth function $f: \mathbb C \to \mathbb R$
and let us study the difference
$\sum_{y \in \mathcal Y_n}
f(\sqrt n \psi(y)) - 
\sum_{y \in \mathcal Y_n}
f(\sqrt n \mathrm d\psi_0(y))$.
Write $\psi(x) = \mathrm d \psi_0 (x) + 
\|x\| r(x)$, where
the limit of $r(x)$ as $x$ goes to zero is zero.
Since $f$ is smooth and compactly supported
it is Lipschitz with 
some Lipschitz constant $C > 0$ and its
support is contained in some disk $\mathcal D$.
Using that $\psi^{-1}$ is Lipschitz near 
the origin we can take a disk 
$\widetilde {\mathcal D}$
such that 
$\psi^{-1}(\mathcal D/\sqrt n) \subset 
\widetilde {\mathcal D}/\sqrt n $
and we may also assume
$\mathrm d \psi_0^{-1}(\mathcal D/\sqrt n) \subset 
\widetilde {\mathcal D}/\sqrt n $.
Consider a compactly supported continuous
function
$g: \mathbb C \to \mathbb R$ satisfying $1_{\widetilde {\mathcal D}} \leq g$. We may write
\begin{align}
\Big|\sum_{y \in \mathcal Y_n}
f(\sqrt n \psi(y)) - 
\sum_{y \in \mathcal Y_n}
f(\sqrt n \mathrm d\psi_0(y))\Big|
&\leq \sum_{y \in \mathcal Y_n}
C \sqrt n|\psi(y) - \mathrm d \psi_0(y)|  
1_{y \in \psi^{-1}(\mathcal D/\sqrt n) 
\cup \mathrm d \psi_0^{-1}(\mathcal D/\sqrt n)}
\nonumber \\
&\leq \sum_{y \in \mathcal Y_n}
C\sqrt n \|y\| |r(y)| 1_{\widetilde 
{\mathcal D}/ \sqrt n}(y)      
\label{eq:SumDelta}    .           
\end{align}
If $R$ is the radius of $\widetilde {\mathcal D}$,
we know that
$\sqrt n\|y\|  1_{\widetilde 
{\mathcal D}/ \sqrt n}(y)   \leq R$.
Since $\lim_{x \to 0} r(x) = 0$,
we know that $r(y) 1_{\widetilde 
{\mathcal D}/ \sqrt n}(y) \leq \varepsilon_n $
for a sequence $\varepsilon_n$ that goes to
zero. Then we can bound
\eqref{eq:SumDelta} by
\[CR\varepsilon_n\sum_{y \in \mathcal Y_n}
g(\sqrt n y) \xrightarrow[n \to \infty]{\dd} 0  .\]
This implies the result.
\end{proof}

\end{document}
